\subsection{Related Work}
\label{subsec:related_work}

Generative AI-based approaches employing diffusion models or flow models enable us to efficiently sample molecular conformations in a multiscale manner. For example, recent structure-generation foundation models, including \AFiii and \Boltzi, employ diffusion-based formulations. In particular, \Boltzii \cite{passaro2025boltz2} introduced \emph{boltz-steering}, a framework that imposes a customized potential to navigate the sampling towards physically plausible conformations. These models largely follow a ``one sequence--one structure'' paradigm, which does not focus on producing diverse conformational ensembles or capturing conformational flexibility. 

Several approaches have been proposed to address this problem. For example, the MSA subsampling technique \cite{delAlamo2022sampling} and its extensions \cite{Wayment-Steele2023MSAClustering, daSilva2024_subsampledAF2, KalakotiWallner2025_AFsample2} increase the diversity of multiple sequence alignments (MSAs) by masking parts of the MSA. 
Another powerful direction is the development of ensemble generative or MD surrogate models. An early contribution in this area is the Boltzmann generator \cite{NoeOlssonKohlerWu2019_BoltzmannGenerators}, which was designed to learn and reproduce the Boltzmann distribution of small, specific peptides. More recently, models have been developed to generate diverse structural ensembles from sequences by learning from extensive MD simulation data across a wide range of proteins \cite{jing2023eigenfold, jing2024alphaflow, Zheng2024_DiG_NatMachIntell, Cheng2024_AlphaFolding, jing2024mdgen, jin2025p2dflow, bioemu2025}. In these methods, obtaining sufficiently long MD trajectories for training requires substantial computational resources, and such models inevitably inherit the force-field biases present in the training data. 

To mitigate this issue, recent studies have explored conditioning ensemble generative models on experimental measurements or physical prior knowledge. Examples include methods that generate structural ensembles conditioned on experimental observables such as X-ray crystallographic electron densities, Nuclear Overhauser effect \cite{maddipatla2025inverse}, structural constraints derived from experimental data \cite{Liu2025_ExEnDiff}, and cryo-EM-consistent modeling \cite{raghu2025multiscale, levy2024solving}. Following this direction, \Model enables AFM-conditioned conformation generation, guiding \AFiii's prior distribution toward modeling more dynamic conformational states captured in HS-AFM images.
